%!TEX root = ../main.tex
\chapter{Related Work}
\label{ch:related}

The work described in this thesis stands at the intersection of four strands
that largely proceed separately: learnt navigation for mobile robots, the
transfer of policies across different bodies, continuous-time networks, and
two-stage training by imitation and reinforcement. Of each, only what is needed
to delimit the contribution is recalled here, distinguishing what is established
from what remains open.

\section{Learnt Navigation for Mobile Robots}
\label{sec:rel-navigation}

Autonomous navigation has long been addressed by composing separately designed
modules --- state estimation, map construction, path planning, control ---
according to the scheme described in Chapter~\ref{ch:background}. The learnt
alternative replaces part of that chain with a policy mapping the observation
directly onto the command, and has in recent years reached performance that
modular composition does not attain.

The most significant results come from flight. A policy trained entirely in
simulation and transferred to real hardware has demonstrated high-speed flight
in unknown natural environments, replacing the whole planning chain with a
network that receives depth and produces a trajectory
\cite{loquercio2021flight}; reinforcement learning has subsequently produced
racing policies competitive with professional human pilots
\cite{kaufmann2023racing}. On the side that bears most directly on this thesis,
continuous-time networks have demonstrated robust aerial navigation outside the
training distribution, preserving the learnt behaviour under environmental
conditions never observed \cite{chahine2023flight}.

These works share an assumption, however: the policy is trained for one
platform. The body, the sensors and the kinematic limits are constants of the
experiment, and not quantities the network observes.

\section{Cross-Embodiment Transfer and Robot Foundation Models}
\label{sec:rel-embodiment}

Independence from the body is not a capability that still awaits demonstration.
The foundational work on datasets aggregated from heterogeneous platforms showed
that training on demonstrations coming from dozens of different robots produces
positive transfer across platforms, with generalist policies that outperform
their single-robot counterparts \cite{openx2024}; on that basis a generation of
vision-language-action models has developed, pursuing generalist control across
different bodies \cite{black2024pi0}. In academic work, morphology-agnostic
control shows that a single controller can generalise to morphologies never
seen, either through transformer representations of the kinematic skeleton
\cite{gupta2022metamorph} or through graph-based policies \cite{luo2025gcnt}.

What characterises such approaches is where independence from the body resides.
In all of them the transfer takes place through the model and the training
process: the abstraction that makes it possible is embedded in the parameters
and fixed before deployment. Two practical consequences follow. The first is
that every new body must be represented in the training data. The second is one
of scale: the models that achieve this generality carry a computational cost
that exceeds what is available on board a small platform, which is measured at
a few inference steps per second even on high-end accelerators
\cite{zhou2026vlaxpu} and collapses by orders of magnitude on genuinely
inexpensive hardware \cite{williams2025litevla}.

The present thesis adopts an intermediate and more modest position: the body is
neither learnt from data nor inferred from the observation, but \emph{declared}
to the network as part of the observation itself, in the form of the kinematic
limits and the model identifier described in
Section~\ref{sec:observation-vector}. The resulting generalisation is limited to
the platforms that descriptor is able to represent, but requires neither
retraining nor additional computational capacity.

\section{Heterogeneous Multi-Robot Policies}
\label{sec:rel-heterogeneous}

It would be incorrect to claim that multi-agent learning assumes homogeneous
agents. It is true that a substantial part of the literature rests on parameter
sharing, with the consequent restriction to the homogeneous case, but
formulations dedicated to learning with heterogeneous agents exist
\cite{zhong2024harl}, as do multi-robot systems with decentralised policies and
explicit heterogeneity, trained in partially observable environments
\cite{bettini2023hetero}.

Particularly close to this work is the approach in which different robots
communicate their own \emph{capabilities}, that is a declarative description of
what each of them is able to do, and the shared policy is conditioned on it
\cite{howell2023capabilities}. The mechanism employed in this thesis is
analogous in form --- an explicit description of the body entering the
observation --- but serves a different purpose: not coordination between agents,
but the reusability of a single policy across a fleet of platforms that do not
cooperate with one another.

\section{Continuous-Time Networks and Liquid Neural Networks}
\label{sec:rel-lnn}

The line of models from which the cells employed in this thesis descend is
recalled in detail in Section~\ref{sec:lnn-background}, and only its terms are
summarised here. Neural ordinary differential equations \cite{chen2018neuralode}
define the hidden state as the solution of a differential equation rather than
as the result of a discrete recurrence; continuous-time recurrent networks
\cite{funahashi1993ctrnn} stabilise their behaviour by introducing a time
constant. Liquid networks make that constant dependent on the input, by explicit
modulation \cite{hasani2021ltc}, by a closed-form approximation of the solution
\cite{hasani2022cfc}, or by making variable the capacitance of the equivalent
circuit as well \cite{farsang2024lrc}. Alongside them stands the sparse wiring
of \emph{Neural Circuit Policies} \cite{lechner2020ncp}, which yields compact
and inspectable control policies.

Two properties make this family pertinent to the problem addressed. The first is
the ability to handle observations at an irregular rate, which in a real robotic
system is the norm rather than the exception. The second is compactness: the
control policies cited employ a number of units orders of magnitude smaller than
traditional recurrent architectures of comparable capability
\cite{lechner2020ncp}, a property that the systematic comparison of the family
confirms on memory footprint as well \cite{zong2025lnnsurvey}.

What the literature does not report, and what this thesis sets out to measure,
is the comparison between the three cells and between the two wirings
\emph{with everything else held equal} on one and the same navigation problem:
the available works each adopt one cell and demonstrate its effectiveness, but
do not place the alternatives on the same bench.

\section{Sparse Encoders for Point Clouds}
\label{sec:rel-encoders}

The representation of the environment the network receives is a cloud of
occupied voxels, and reducing it to a descriptor of fixed length is the task of
the encoder. The two available families are the one operating directly on the
unordered set of points \cite{qi2017pointnet} and the one that discretises space
on a grid and applies convolutions to it. In the latter, \emph{submanifold}
convolutions keep the computation anchored to the occupied sites alone instead
of densifying the grid \cite{graham2018submanifold}, a property that makes
tractable an input in which almost all cells are empty. It is the solution
adopted in Section~\ref{sec:sparse-encoder}, and the reason for the choice is
that the local map produced by the pipeline is already quantised on a grid:
operating on continuous points would mean disregarding a structure the upstream
system has already imposed.

\section{Imitation Followed by Reinforcement}
\label{sec:rel-imitation-rl}

Imitation learning makes it possible to obtain an initial behaviour from
recorded demonstrations, but suffers from a structural defect: the learnt policy
is evaluated on the distribution of the states visited by the demonstrator,
whereas in execution it visits the states produced by its own actions, and the
error accumulates along the trajectory \cite{ross2011dagger}. Reinforcement
learning does not have that defect, since it learns on the distribution it
generates itself, but requires a number of interactions that a real robot cannot
provide.

The combination of the two --- imitation to obtain a starting behaviour,
reinforcement to refine it on the correct distribution --- is the strategy
adopted in Chapter~\ref{ch:training}. The reinforcement stage employs a
recurrent variant of proximal policy optimisation \cite{schulman2017ppo} with
generalised advantage estimation \cite{schulman2016gae}, and the design of the
reward abides by the principle that a shaping term must not alter the optimal
policy \cite{ng1999policy}. The implementation choices that in such algorithms
bear on the result as much as the formulation itself are documented in
\cite{engstrom2020implementation}, and Section~\ref{sec:training-rl} states
explicitly those adopted here.

\section{Positioning of This Work}
\label{sec:rel-positioning}

With respect to the strands recalled above, the contribution of this thesis is
circumscribed and can be stated by difference. It does not propose a new
continuous-time architecture, but compares the three existing ones under equal
conditions on a real navigation problem. It does not pursue independence from
the body through training on many bodies, but obtains it by declaring the body
in the observation, at the cost of a generalisation limited to what the
descriptor represents. It does not address coordination between agents, since
the fleet employed in reinforcement training serves to produce experience in
parallel and not to cooperate. And it takes as a constraint, rather than as a
subsequent question, that the resulting policy must run on the on-board hardware
of the platform that employs it.
